% Diagrama de classes do domínio do sistema de agendamento (página em pé).
% Fonte: models.py das apps accounts, schedule, bookings e notices.
% Classes de apoio (importação, auditoria, fila de e-mails) foram omitidas.
\begin{tikzpicture}[
    font=\scriptsize,
    cls/.style={draw, rectangle split, rectangle split parts=3,
                text width=3.4cm, align=left, anchor=north, inner sep=2.5pt,
                rectangle split part fill={black!8,white,white}},
    cls2/.style={draw, rectangle split, rectangle split parts=2,
                text width=3.4cm, align=left, anchor=north, inner sep=2.5pt,
                rectangle split part fill={black!8,white}},
    m/.style={font=\tiny, inner sep=1.2pt},
    rl/.style={font=\tiny\itshape, inner sep=1.2pt, fill=white},
    comp/.style={{Diamond[length=2.6mm, width=1.8mm]}-},
]
\newcommand{\cname}[1]{\centering\textbf{\footnotesize #1}\par}
\def\ca{1.85}\def\cb{6.85}\def\cc{11.45}\def\cd{16.05}
\def\ra{21.0}\def\rb{18.4}\def\rc{15.2}\def\rd{9.8}\def\re{6.9}

% ---------------- coluna 1: acesso ----------------
\node[cls2] (permissao) at (\ca,\ra) {\cname{Permissão}
  \nodepart{second} - codigo: texto\\- descricao: texto};
\node[cls2] (perfil) at (\ca,\rb) {\cname{Perfil}
  \nodepart{second} - codigo: texto\\- nome: texto\\- descricao: texto};
\node[cls] (usuario) at (\ca,\rc) {\cname{Usuário}
  \nodepart{second} - nome: texto\\- email: texto\\- nome\_usuario: texto\\- ativo: lógico\\- superadmin: lógico
  \nodepart{third} + tem\_permissao(cod)};
\node[cls2] (notificacao) at (\ca,\rd) {\cname{Notificação}
  \nodepart{second} - titulo: texto\\- mensagem: texto\\- lida\_em: data-hora};

% ---------------- coluna 2: agenda ----------------
\node[cls2] (outronome) at (\cb,\ra) {\cname{Outro nome}
  \nodepart{second} - nome: texto};
\node[cls2] (agendamento) at (\cb,\rb) {\cname{Agendamento}
  \nodepart{second} - inicio: data-hora\\- fim: data-hora\\- descricao: texto\\- qtd\_alunos: inteiro};
\node[cls] (solicitacao) at (\cb,\rc) {\cname{Solicitação}
  \nodepart{second} - data: data\\- inicio: hora\\- fim: hora\\- descricao: texto\\- qtd\_alunos: inteiro\\- situacao: texto
  \nodepart{third} + submeter()\\+ aprovar(analista)\\+ rejeitar(motivo)\\+ sugerir(alternativa)\\+ cancelar(motivo)};
\node[cls2] (vinculo) at (\cb,\rd) {\cname{Vínculo docente}
  \nodepart{second} - atual: lógico};
\node[cls2] (docente) at (\cb,\re) {\cname{Docente}
  \nodepart{second} - nome: texto};

% ---------------- coluna 3: espaços e grade ----------------
\node[cls2] (espaco) at (\cc,\rb) {\cname{Espaço}
  \nodepart{second} - nome: texto\\- origem: texto\\- ativo: lógico};
\node[cls2] (ensalamento) at (\cc,\rc) {\cname{Ensalamento}
  \nodepart{second} - dia\_semana: texto\\- origem: texto\\- vigente\_de: data\\- vigente\_ate: data};
\node[cls2] (disciplina) at (\cc,\rd) {\cname{Disciplina}
  \nodepart{second} - nome: texto};
\node[cls2] (oferta) at (\cc,\re) {\cname{Oferta de disciplina}
  \nodepart{second} \vphantom{x}};

% ---------------- coluna 4 ----------------
\node[cls2] (faixa) at (\cd,\ra) {\cname{Faixa de horário}
  \nodepart{second} - inicio: hora\\- fim: hora};
\node[cls2] (semestre) at (\cd,\rb) {\cname{Semestre}
  \nodepart{second} - codigo: texto\\- inicio: data\\- fim: data};
\node[cls2] (ocorrencia) at (\cd,\rc) {\cname{Ocorrência cancelada}
  \nodepart{second} - data: data\\- tipo: texto};
\node[cls2] (turma) at (\cd,\rd) {\cname{Turma}
  \nodepart{second} - rotulo: texto};
\node[cls2] (curso) at (\cd,\re) {\cname{Curso}
  \nodepart{second} - nome: texto\\- etapa: texto};

% ================= associações =================
% --- acesso (verticais na coluna 1)
\draw (perfil.north) -- node[rl, right] {concede} (permissao.south)
      node[m, pos=0.12, left] {*} node[m, pos=0.88, left] {*};
\draw (usuario.north) -- node[rl, right] {tem} (perfil.south)
      node[m, pos=0.12, left] {*} node[m, pos=0.88, left] {0..1};
\draw (usuario.south) -- node[rl, right] {recebe} (notificacao.north)
      node[m, pos=0.08, left] {1} node[m, pos=0.92, left] {*};

% --- usuário e solicitação (horizontais)
\path (usuario.north east) ++(0,-0.75) coordinate (s1);
\path (usuario.north east) ++(0,-1.75) coordinate (s2);
\draw (s1) -- node[rl, above] {solicita} (solicitacao.west |- s1)
      node[m, pos=0.1, below] {1} node[m, pos=0.9, below] {*};
\draw (s2) -- node[rl, above] {analisa} (solicitacao.west |- s2)
      node[m, pos=0.12, below] {0..1} node[m, pos=0.9, below] {*};

% --- agenda (verticais na coluna 2)
\draw (solicitacao.north) -- node[rl, right] {gera} (agendamento.south)
      node[m, pos=0.12, left] {1} node[m, pos=0.88, left] {0..1};
\draw (solicitacao.south) -- node[rl, right] {justifica} (vinculo.north)
      node[m, pos=0.08, left] {*} node[m, pos=0.92, left] {0..1};
\draw (docente.north) -- node[rl, right] {assume} (vinculo.south)
      node[m, pos=0.12, left] {1} node[m, pos=0.88, left] {*};

% --- espaço
\path (agendamento.north east) ++(0,-0.9) coordinate (a1);
\draw (a1) -- node[rl, above] {ocupa} (espaco.west |- a1)
      node[m, pos=0.1, below] {*} node[m, pos=0.9, below] {1};
\draw (solicitacao.north east) -- node[rl, sloped, below] {é pedida em} (espaco.south west)
      node[m, pos=0.1, above left] {*} node[m, pos=0.85, above left] {1};
\draw[comp] (espaco.north west) -- (outronome.south east)
      node[m, pos=0.85, below right] {*};

% --- ensalamento
\draw (ensalamento.north) -- node[rl, right] {usa} (espaco.south)
      node[m, pos=0.12, left] {*} node[m, pos=0.88, left] {0..1};
\draw (ensalamento.south) -- (disciplina.north)
      node[m, pos=0.06, left] {*} node[m, pos=0.94, left] {1};
\path (ensalamento.north east) ++(0,-0.9) coordinate (e1);
\draw[comp] (e1) -- (ocorrencia.west |- e1)
      node[m, pos=0.85, above] {*};
\draw ([yshift=-3mm]ensalamento.north east) -- (semestre.south west)
      node[m, pos=0.15, below right] {*} node[m, pos=0.85, below right] {1};
\draw (ensalamento.north east) -- (faixa.south west)
      node[m, pos=0.06, left] {*} node[m, pos=0.95, left] {1};
\draw (ensalamento.south east) -- (turma.north west)
      node[m, pos=0.12, below left] {*} node[m, pos=0.88, below left] {1};
\draw (docente.east) -- (\cb+2.3,\re-0.5) -- node[rl, sloped, above, pos=0.55] {ministra} (\cb+2.3,\rc-1.75)
      -- (ensalamento.south west |- 0,\rc-1.75);
\node[m, above right] at (docente.east) {*};
\node[m, above left] at (ensalamento.south west |- 0,\rc-1.75) {*};

% --- grade
\draw (semestre.north) -- node[rl, right] {define} (faixa.south)
      node[m, pos=0.12, left] {1} node[m, pos=0.88, left] {*};
\draw (vinculo.south east) -- (oferta.north west)
      node[m, pos=0.12, above right] {*} node[m, pos=0.85, above right] {1};
\draw (oferta.north) -- (disciplina.south)
      node[m, pos=0.12, left] {*} node[m, pos=0.88, left] {1};
\draw (oferta.north east) -- (turma.south west)
      node[m, pos=0.12, above left] {*} node[m, pos=0.88, above left] {1};
\draw (curso.north) -- node[rl, right] {possui} (turma.south)
      node[m, pos=0.12, left] {1} node[m, pos=0.88, left] {*};
\draw (oferta.south) -- (\cc,\re-1.6) -- (\cd+2.3,\re-1.6) -- (\cd+2.3,\rb-0.8) -- (semestre.east |- 0,\rb-0.8)
      node[m, pos=0.9, above] {1};
\node[m, below left] at (oferta.south) {*};
\end{tikzpicture}
